%!TEX program = xelatex
\documentclass[10.5pt]{article}

\usepackage[UTF8]{ctex}
\usepackage[margin=1in]{geometry}
\usepackage{amsmath,amssymb}
\usepackage{booktabs}
\usepackage{tabularx}
\usepackage{enumitem}
\usepackage{natbib}
\usepackage[colorlinks=true,linkcolor=blue,citecolor=blue,urlcolor=blue]{hyperref}
\usepackage{microtype}

\setlist{nosep,leftmargin=*}
\setlength{\parindent}{2em}
\setlength{\parskip}{0.35em}
\setlength{\emergencystretch}{3em}
\newcommand{\code}[1]{\texttt{#1}}

\title{市场状态自适应的多智能体组合管理\\
\large 文献综述与研究设计草案}
\author{}
\date{2026年9月7日}

\begin{document}
\maketitle

\begin{abstract}
本文综述多智能体强化学习、LLM金融智能体、智能体市场仿真与量化组合管理的相关研究，并据此提出一个可复现的增量式研究设计。研究拟以已有的MASA和FinRL框架为基础，引入市场状态识别、基于置信度的智能体动态加权以及包含换手率与尾部风险惩罚的奖励函数。模型改动保持在可解释且易于消融的范围内，主要比较原始模型、静态加权集成、单智能体强化学习和传统组合基线。本文目前仅为文献综述与研究设计，不报告尚未完成的实证结果。
\end{abstract}

\section{研究背景与问题界定}

深度强化学习已经被用于将组合管理表述为连续决策问题，典型研究以价格、技术指标和持仓状态构造观察变量，并以组合收益或风险调整收益定义奖励函数\citep{JiangXuLiang2017,LiuYangChen2020}。与单一策略网络相比，多智能体框架可以把不同的投资信号、信息来源或风险偏好拆分为相对独立的决策单元，再通过协调机制生成最终交易动作。近年的LLM金融智能体进一步把基本面、新闻、情绪、技术分析和风险管理等角色显式化\citep{YuLiChen2023,XiaoSunLuo2024,YuYaoLi2024}。

然而，现有工作之间的研究对象并不完全相同。数值型多智能体强化学习主要关注连续控制、策略协作和组合风险；LLM智能体研究则更强调自然语言信息处理、工具调用和角色化推理；智能体市场仿真研究关注投资者异质性、订单流和价格形成机制\citep{ByrdHybinetteBalch2019,LopezLira2025,LiLiuLiu2024,YangZhangWu2025}。因此，本文将“多智能体”限定为一组能够产生可记录交易动作或风险信号、并通过显式聚合机制影响组合决策的决策单元，而不把普通的多模块软件流水线自动等同于多智能体模型。

本文拟回答的核心问题是：
\begin{quote}
在市场状态发生变化时，基于市场状态、智能体置信度和尾部风险的动态权重机制，能否相对于静态集成和单智能体基线提高组合管理的样本外风险调整表现？
\end{quote}

这一问题允许采用已有开源实现作为基线，同时把贡献集中在一个明确的聚合机制和严格的样本外比较上，而不是把代码规模或提示词数量本身作为创新。

\section{多智能体强化学习与组合管理}

早期深度强化学习组合管理研究通常使用一个策略网络直接输出资产权重或交易动作。EIIE框架将资产价格和技术指标输入策略网络，并以组合收益驱动端到端学习\citep{JiangXuLiang2017}；FinRL则提供了包含环境、算法和评估流程的开源库，降低了量化强化学习实验的实现成本\citep{LiuYangChen2020}。这些研究为动作空间、交易成本和基准比较提供了重要起点，但单一策略网络将不同类型的信号压缩在同一表示中，难以观察不同决策来源在市场状态变化中的相对作用。

多智能体强化学习提供了另一种分解方式。MASA框架把多智能体、自适应机制和深度强化学习结合起来，用于动态组合风险管理\citep{LiTamYeung2024}。从研究设计角度看，这类模型的关键不只是智能体数量，而是智能体之间如何分工、如何共享信息以及如何将局部策略转换为全局组合。若聚合权重固定，增加智能体可能只带来参数规模和计算成本的上升；若聚合权重随市场状态变化，则可以检验不同策略在趋势、震荡和高波动状态下是否具有不同的边际价值。

本文计划采用一个较小的结构扩展：保留各个基础智能体的局部策略网络，在其上增加一个门控模块。门控模块以市场状态向量、智能体近期风险调整表现、输出置信度和预测风险为输入，生成非负且和为一的动态权重。最终动作写为
\begin{equation}
    a_t=\sum_{i=1}^{N} w_{i,t}a_{i,t},\qquad
    w_{i,t}=\frac{\exp(g_i(z_t,c_{i,t},q_{i,t}))}
    {\sum_{j=1}^{N}\exp(g_j(z_t,c_{j,t},q_{j,t}))},
\end{equation}
其中，$z_t$表示市场状态，$c_{i,t}$表示第$i$个智能体的置信度，$q_{i,t}$表示其风险估计，$g_i(\cdot)$为门控网络。该设定并不预设某一智能体在所有时期都占优，而是允许权重随状态变化。

\section{LLM金融智能体与角色化交易框架}

LLM金融智能体研究大致形成三条路线。第一条路线将基本面、技术面、新闻和情绪分析分配给不同角色，再由交易员或投资组合经理汇总结果。TradingAgents将分析师团队、多空研究员、交易员和风险管理团队组织为协作图\citep{XiaoSunLuo2024}；ContestTrade进一步把数据压缩、研究提案和内部竞赛结合起来\citep{ZhaoSunJiang2025}。第二条路线强调记忆、人格和语言反馈，例如FinMem使用分层记忆和角色设计改善交易智能体的连续决策\citep{YuLiChen2023}，FinCon则通过经理—分析师结构和概念性语言强化组织多智能体投资决策\citep{YuYaoLi2024}。第三条路线使用多模态或工具增强的基础智能体处理价格、图表、新闻和其他外部资料\citep{ZhangZhaoXia2024}。

上述研究证明了角色化设计的可行性，但对组合模型的性能比较仍面临三个问题。其一，文本推理、数值决策和执行模拟经常被放在同一系统中，导致难以判断收益变化来自信息处理、策略聚合还是数据与提示词差异。其二，许多实验重点放在累计收益或单一回测曲线上，对换手、滑点、容量、尾部损失和模型调用成本的处理不够统一。其三，LLM输出具有随机性，模型版本、温度、上下文和工具返回结果都会影响交易动作。因此，本文将LLM框架作为可选扩展基线，而把主模型先限定为数值型、可重复的多智能体强化学习系统。

\section{智能体市场仿真与金融机制}

如果研究目标是解释市场价格如何由异质投资者形成，则市场仿真比单纯的组合回测更合适。ABIDES提供了面向AI研究的高保真市场仿真环境，可以表示限价订单簿、交易所和不同类型的交易智能体\citep{ByrdHybinetteBalch2019}。MarS和TwinMarket分别探索了生成式基础模型驱动的金融市场仿真与行为、社会信号模拟\citep{LiLiuLiu2024,YangZhangWu2025}。Lopez-Lira的研究则直接把LLM智能体放入市场模拟中，用于检验金融理论，而不是只比较一个策略的净值\citep{LopezLira2025}。

这类研究对本文有两点启发。第一，智能体的异质性应当对应可观察的投资信号、风险偏好或信息集，而不只是不同的提示词。第二，模型评价不应只看收益，还应检验价差、成交量、波动聚集、价格冲击和市场效率等对象。本文第一阶段仍以组合管理为主，以保证研究规模可控；若基础模型显示出稳定的状态依赖性，第二阶段可将动态聚合模块嵌入订单簿仿真，考察其对交易拥挤和流动性的影响。

\section{代码先于论文的开源项目及其研究用途}

当前GitHub上已经出现一些代码较完整、但在本次检索范围内尚未检出对应论文或预印本的项目。这里的“未检出”只描述检索状态，不构成不存在论文的证明。

\begin{table}[htbp]
\centering
\caption{可作为实验起点的开源项目}
\label{tab:opensource}
\small
\begin{tabularx}{\textwidth}{lXl}
\toprule
项目 & 主要功能与研究用途 & 论文状态（截至2026-09-07）\\
\midrule
\href{https://github.com/virattt/ai-hedge-fund}{ai-hedge-fund} & 多角色投研、风险管理和回测；适合作为角色化集成的工程基线 & 暂未检出对应论文；README定位为概念验证 \\
\href{https://github.com/QuantSingularity/LLM-Powered-Multi-Agent-Frameworks-for-Algorithmic-Trading}{QuantSingularity} & 分析、决策、风控、执行与解释五类角色框架 & 暂未检出对应论文 \\
\href{https://github.com/harrypapadakis/StockSim}{StockSim} & 订单簿和异步多智能体市场仿真；适合研究执行和价格形成 & 暂未检出对应论文 \\
\href{https://github.com/arteemg/AutoHypothesis}{AutoHypothesis} & 自动提出、回测和滚动验证量化假设；适合作为研究流程基线 & 暂未检出对应论文，但与近期自动化alpha研究框架存在重叠 \\
\href{https://github.com/HKUDS/Vibe-Trading}{Vibe-Trading} & agentic量化研究、因子分析、回测和数据工具 & 暂未检出对应论文；更接近研究平台 \\
\bottomrule
\end{tabularx}
\end{table}

相较之下，TradingAgents、ContestTrade和MASA均已有公开论文，分别对应角色化LLM交易、内部竞赛机制和多智能体自适应组合风险管理\citep{XiaoSunLuo2024,ZhaoSunJiang2025,LiTamYeung2024}。因此，本文将这些项目视为可引用的模型先例和基线来源，而不是声称重新发现一个尚未被研究的方向。

\section{拟议模型与可检验假设}

拟议模型暂命名为Regime-Adaptive Confidence-Weighted Multi-Agent Portfolio（RAC-MAP）。模型由五类基础智能体组成：动量、价值、情绪、波动率和宏观状态。基础智能体可以采用相同的PPO或SAC算法，也可以采用不同的观察变量；门控模块负责把局部动作转换为组合动作。奖励函数为
\begin{equation}
    R_t=r_t-\lambda\,Turnover_t-\eta\,CVaR_t,
\end{equation}
其中$r_t$为组合收益，$Turnover_t$为调仓换手率，$CVaR_t$为滚动窗口内的条件风险价值。该改动有三个可单独检验的组成部分：市场状态门控、置信度加权和风险成本。

据此形成以下可检验假设：
\begin{enumerate}
    \item 在市场状态转换较频繁的区间，动态门控的样本外风险调整表现优于固定权重集成。
    \item 在预测置信度与实际误差具有信息含量时，置信度加权能够降低组合回撤和尾部损失。
    \item 将换手和尾部风险纳入奖励函数后，模型的成本调整表现和策略稳定性优于只最大化组合收益的模型。
\end{enumerate}

这些假设是模型层面的可检验命题，不能在实验完成前写成已经成立的结论。

\section{数据来源与最新性安排}

\subsection{数据源}

正式实验优先使用具有明确历史版本和点时信息的数据。CRSP/Compustat和WRDS（\url{https://wrds-www.wharton.upenn.edu/}）适合美国股票收益、财务报表和退市信息，但通常需要机构许可；SEC EDGAR的提交记录和Company Facts接口（\url{https://www.sec.gov/edgar/sec-api-documentation}）可提供免费、可追溯的公司披露信息。FRED（\url{https://fred.stlouisfed.org/}）可用于宏观和波动率序列，Kenneth French Data Library（\url{https://mba.tuck.dartmouth.edu/pages/faculty/ken.french/data_library.html}）可用于因子控制。Yahoo Finance的chart接口（\url{https://query1.finance.yahoo.com/v8/finance/chart/SPY?range=10y&interval=1d&events=history}）便于原型和公开复现，但其数据授权、修订和退市覆盖不宜直接等同于机构级数据库。

截至2026年9月7日的可用性核验中，Yahoo Finance公开接口返回SPY日频数据至2026年9月4日、BTC-USD日频数据至2026年9月7日；SEC EDGAR的公司申报接口（例如Apple的提交记录：\url{https://data.sec.gov/submissions/CIK0000320193.json}）可以直接访问。上述日期是接口可见的最新观测，不是已经冻结的研究样本。正式样本建议以2025年12月31日为最近一个完整年度终点，2026年数据仅用于后续滚动更新或外部监测。

\begin{table}[htbp]
\centering
\caption{拟用数据源与用途}
\label{tab:data}
\small
\begin{tabularx}{\textwidth}{lXll}
\toprule
数据源 & 主要变量与用途 & 时效性 & 研究阶段\\
\midrule
CRSP/Compustat（WRDS） & 股票收益、财务报表、退市和公司特征；用于正式资产定价检验 & 依机构订阅更新 & 主论文优先 \\
SEC EDGAR & 点时财务披露、申报日期和文本信息；用于情绪或信息处理变量 & 持续更新 & 主论文/扩展 \\
FRED & 宏观利率、波动率及金融条件变量 & 持续更新 & 控制变量与状态识别 \\
Kenneth French Data Library & 市场、规模、价值、盈利和投资等因子 & 定期更新 & 风险调整评价 \\
Yahoo Finance chart API & 公开日频价格；便于原型、复现和加密资产扩展 & 接口实时返回；需冻结下载日期 & 原型与敏感性分析 \\
\bottomrule
\end{tabularx}
\end{table}

\subsection{样本外评估协议}

主实验采用按时间排序的训练、验证和测试区间，不随机打乱观测。每个实验版本固定数据快照、代码提交、随机种子、交易成本参数和模型超参数。测试指标包括年化收益、年化波动率、Sharpe比率、Sortino比率、最大回撤、Calmar比率、CVaR、换手率和成本调整后的收益。模型比较还应报告不同随机种子、不同资产池和不同市场状态下的分布，而非只报告一个最优回测结果。

基线至少包括买入并持有、等权组合、风险平价、单智能体PPO/SAC、原始FinRL或MASA、静态权重多智能体，以及RAC-MAP的三组消融版本。若加入LLM基线，应固定模型版本、提示词、温度、外部数据快照和调用预算，并把文本推理成本单独记录。

\section{文献评述与预期贡献边界}

现有研究已经覆盖了金融LLM智能体、深度强化学习组合管理和智能体市场仿真的主要技术路线\citep{DingLiWang2024,XiaYouWang2026}。本文不把“首次使用多智能体”作为贡献，也不把单次回测中的收益提升直接解释为金融机制。预期贡献限定为三点：第一，在可复现的数值型多智能体组合框架中引入市场状态、置信度和尾部风险的统一聚合机制；第二，使用静态集成、单智能体和传统组合模型构成可解释的基线矩阵；第三，以多时期、多随机种子、交易成本和风险指标检验模型改动是否具有稳定的样本外价值。

如果动态门控仅在单一市场、单一时期或无成本回测中有效，则其结论只能被表述为有限的模型比较结果。若它在多个市场、多个状态和不同成本假设下均保持优势，并且消融实验能够显示优势来自预设的门控或风险机制，论文才具备进一步冲击金融量化和金融科技期刊的基础。

\bibliographystyle{plainnat}
\bibliography{literature_review_zh}

\end{document}
